% !PRINTED 521
\weberpage{521}
Im ersten Falle kommen nach § 105 nur die beiden Kongruenzen
\[
A \equiv 1, \quad A \equiv 5\ (\mathrm{mod}\ 8),
\]
im zweiten Falle nur die beiden Kongruenzen
\[
A \equiv 1, \quad A \equiv 3\ (\mathrm{mod}\ 8)
\]
vor, und zwar wieder in gleich viel Formenklassen. Die vollständige Zerlegung der Klassengleichung nach den Geschlechtern erfordert, wenn $m$ durch 8 teilbar ist, immer die Adjunktion der Wurzeln aus \gesperrt{sämtlichen} in $m$ aufgehenden Primzahlen, einschließlich 2.

Man bemerkt, daß in diesen Betrachtungen ein neuer Beweis der von \gesperrt{Gauss} und \gesperrt{Dirichlet} bewiesenen Sätze über die Existenz der Geschlechter enthalten ist, freilich nur für negative Diskriminanten.

\begin{center}
\S\ 139. \gesperrt{Beispiele.}
\end{center}

Zur wirklichen Ausrechnung der Zerfällung der Klassengleichung sind die Formeln des vorigen Paragraphen nur in beschränktem Maße anwendbar wegen des zu hohen Grades der Transformationsgleichungen. Wir werden nachher in einem Falle eine wenigstens nahe verwandte Methode zur Anwendung bringen. Ist die Klassengleichung bekannt, so läßt sich meist leichter die Zerlegung direkt finden, indem man einen Ansatz von bekannter Form mit unbestimmten Koeffizienten macht; so findet man aus den Formeln § 130, (2), (29), (37) die folgenden Teilgleichungen, worin das positive Zeichen der Quadratwurzel dem Hauptgeschlecht entspricht:
\[
m = 50, \quad f_1\!\left(\sqrt{-50}\right) = \sqrt[4]{2}\, x,
\]
\[
x^3 - x^2 = \frac{1+\sqrt5}{2}(x+1),
\]
\[
m = 26, \quad f_1\!\left(\sqrt{-26}\right)^2 = \sqrt2\, x,
\]
\[
x^3 - x^2 = \frac{3+\sqrt{13}}{2}(x+1),
\]
\[
m = 41, \quad z = \frac{f\!\left(\sqrt{-41}\right)^2}{\sqrt2} + \frac{\sqrt2}{f\!\left(\sqrt{-41}\right)^2},
\]
\[
z^2 - \frac{5+\sqrt{41}}{2} z + \frac{7+\sqrt{41}}{2} = 0.
\]
